\documentclass[12pt,border=0pt]{standalone}
\input{style.tex}
\usepackage{mathrsfs}
% Figure 1 of manuscript 14, version 14.112: frames 239.65 pt high (were 247 pt), the least reduction that lets
% the figure and its shortened caption fit the text height (LaTeX had reported the float 14.65 pt too large). (Whole page, two panels stacked vertically, exactly 16.0 cm wide with the
% tick labels, inside the text block): the Langevin magnetisation curve of the Eulerian distribution,
% m(H) = 1/2 + L(H/2)/2 (data: make_fig01_langevin_page.py).
% (a) and (b) have identical frames. (a) shows -10 <= H <= 10 (author's instruction of 8 November 2026); the dashed
%     rectangle in (a) is the window [H_0-2.5,H_0+2.5] x [m(H_0)-0.1375,m(H_0)+0.1375] of (b), centred at z_0
%     (H_0 = -1/sqrt2, m(H_0) = 0.441559707648; v14.117), with the same shape as the frame of (b): its width fraction
%     5/20 equals its height fraction 0.275/1.1. In (b) the curve meets the bottom of the frame at H = -2.6052744785
%     and the top at H = 0.9632926514 (data: make_fig01_langevin_page.py).
% (a) also the inflectional tangent m = 1/2 + H/12 for -5.2 <= H <= 5.2 (30% longer than the former -4 <= H <= 4)
%     and the inflection point (0, 1/2).
% (b) the window of half-width 2.5 centred at z_0 (author's instructions of 8 November 2026) with the fields of z_0
%     and z_1 = S z_0.
% In both panels the inflection point is a white disc with a black rim drawn on top: no line crosses it.
% y axes labelled m(H).
% Typography as Figure 2: frame and ticks 0.3 pt, drawn on top; dotted guides 0.6 pt, dashed guides 0.7 pt.
\pgfplotsset{fig1frame/.style={axis on top=true,axis line style={line width=0.3pt,black},every tick/.style={line width=0.3pt,black}}}
\newcommand{\InsetLabel}[1]{{\fontsize{15}{15}\selectfont #1}}
\begin{document}
\begin{tikzpicture}
\begin{axis}[ojac,fig1frame,scale only axis=true,width=401.26pt,height=239.65pt,name=top,
 xmin=-10,xmax=10,ymin=-0.05,ymax=1.05,
 xlabel={$H=\log u$},
 ylabel={$m(H)$},ylabel shift=-3pt,
 xtick={-10,-8,-6,-4,-2,0,2,4,6,8,10},ytick={0,0.25,0.5,0.75,1},
 yticklabel style={/pgf/number format/fixed,/pgf/number format/precision=2},
 clip=false]
  \addplot[black,dotted,line width=0.6pt,no marks] coordinates {(-10,0) (10,0)};
  \addplot[black,dotted,line width=0.6pt,no marks] coordinates {(-10,1) (10,1)};
  \addplot[black,dotted,line width=0.6pt,no marks] coordinates {(-10,0.5) (10,0.5)};
  % window of panel (b), same aspect ratio as the frame of (b)
  \draw[black!65,dashed,line width=0.7pt] (axis cs:-3.207106781187,0.304059707648) rectangle (axis cs:1.792893218813,0.579059707648);
  \node[anchor=north west,black!75,inner sep=0.8pt] at (axis cs:1.792893218813,0.304059707648) {\InsetLabel{(b)}};
  \addplot[figthreeblue,line width=1.1pt,no marks,restrict x to domain=-10:10] table[x=H,y=m] {../data/fig01_langevin_S.dat};
  % the inflectional tangent m = 1/2 + H/12, centred at the inflection point (0, 1/2)
  \addplot[black,line width=0.4pt,no marks] coordinates {(-5.2,0.06666667) (5.2,0.93333333)};
  % the inflection point, as in (b): a white disc drawn on top, so that no line crosses it
  \addplot[only marks,mark=*,mark size=2.6pt,mark options={fill=white,draw=black,line width=0.6pt}] coordinates {(0,0.5)};
  \node[anchor=north west,inner sep=2pt] at (rel axis cs:0.012,0.93) {\InsetLabel{(a)}};
  \node[font=\small,anchor=north west] at (axis cs:-7.08,0.985) {$m=1$};
  \node[font=\small,anchor=south east] at (axis cs:9.77,0.015) {$m=0$};
\end{axis}
\begin{axis}[ojac,fig1frame,scale only axis=true,width=401.26pt,height=239.65pt,
 at={($(top.south west)+(0,-1.75cm)$)},anchor=north west,
 xmin=-3.207106781187,xmax=1.792893218813,ymin=0.304059707648,ymax=0.579059707648,
 xlabel={$H=\log u$},
 ylabel={$m(H)$},ylabel shift=-3pt,
 xtick={-3,-2.5,-2,-1.5,-1,-0.5,0,0.5,1,1.5},ytick={0.35,0.4,0.45,0.5,0.55},
 xticklabel style={/pgf/number format/fixed,/pgf/number format/precision=1},
 yticklabel style={/pgf/number format/fixed,/pgf/number format/precision=2},
 clip=false]
  \addplot[figthreeblue,line width=1.1pt,no marks] table[x=H,y=m] {../data/fig01_langevin_zoom.dat};
  \addplot[black,dotted,line width=0.6pt,no marks] coordinates {(-3.207106781187,0.5) (1.792893218813,0.5)};
  % guide lines
  \addplot[figthreeblue,dashed,line width=0.7pt,no marks] coordinates {(-0.707106781187,0.304059707648) (-0.707106781187,0.441559707648) (-3.207106781187,0.441559707648)};
  \addplot[ojacred,dashed,line width=0.7pt,no marks] coordinates {(-1.414213562373,0.304059707648) (-1.414213562373,0.385899079161) (-3.207106781187,0.385899079161)};
  % the two points and the inflection point
  \addplot[only marks,mark=*,mark size=3pt,figthreeblue] coordinates {(-0.707106781187,0.441559707648)};
  \addplot[only marks,mark=*,mark size=3pt,ojacred] coordinates {(-1.414213562373,0.385899079161)};
  % arrow from the blue point to the red point: the modular inversion S as doubling of the field
  \draw[-{Latex[length=6pt,width=4.5pt]},black,line width=0.7pt,shorten >=4pt,shorten <=4pt]
    (axis cs:-0.707106781187,0.441559707648) to[bend left=28] (axis cs:-1.414213562373,0.385899079161);
  \node[font=\small,anchor=south] at (axis cs:-0.97,0.451) {$S$};
  % labels of the points
  \node[font=\small,anchor=west,figthreeblue] at (axis cs:-0.62,0.428) {$z_0$};
  \node[font=\small,anchor=north west,ojacred] at (axis cs:-1.37,0.380) {$z_1=Sz_0$};
  \node[font=\small,anchor=south west,figthreeblue] at (axis cs:-3.177,0.444) {$m(H_0)=0.4416$};
  \node[font=\small,anchor=south west,ojacred] at (axis cs:-3.177,0.3915) {$m(H_1)=0.3859$};
  \node[font=\small,anchor=south east] at (axis cs:1.753,0.502) {$m=\dfrac12$};
  \node[anchor=north west,inner sep=2pt] at (rel axis cs:0.012,0.985) {\InsetLabel{(b)}};
  % the inflection point: a white disc drawn last, so that no line crosses it
  \addplot[only marks,mark=*,mark size=2.6pt,mark options={fill=white,draw=black,line width=0.6pt}] coordinates {(0,0.5)};
\end{axis}
\end{tikzpicture}
\end{document}
